\chapter{Praktikum 3: SPNL dan Pencocokan Kurva}
\CP{Mahasiswa mampu menurunkan dan menjalankan Newton untuk beberapa variabel, merumuskan OLS linear sebagai SPL, membentuk Jacobian model nonlinear, menjalankan Gauss--Newton, dan mengevaluasi hasil dengan galat serta grafik konvergensi.}
\Target{Dua bagian terpisah: SPNL dan pencocokan kurva. Kode terdiri atas fungsi algoritma yang dapat digunakan ulang dan \emph{driver} berisi data, parameter, tabel, serta visualisasi.}

\section{Bagian A --- Newton untuk Sistem Persamaan Nonlinear}
\subsection{Dari akar skalar ke beberapa variabel}
Pada satu variabel, Newton mencari akar $f(x)=0$ dengan menggunakan turunan $f'(x)$. Pada $n$ variabel, kita mencari vektor $\mathbf{x}=(x_1,\ldots,x_n)^T$ yang membuat \emph{seluruh} persamaan bernilai nol:
\begin{equation}
\mathbf{F}(\mathbf{x})=\begin{bmatrix}f_1(\mathbf{x})\\ \vdots\\ f_n(\mathbf{x})\end{bmatrix}=\mathbf{0}.
\end{equation}
Turunan $f'(x)$ digantikan oleh \textbf{matriks Jacobian}, yaitu kumpulan turunan parsial:
\begin{equation}
J_F(\mathbf{x})=\left[\frac{\partial f_i}{\partial x_j}(\mathbf{x})\right]_{i,j=1}^{n}.
\end{equation}
Jangan mencampurkan nama \emph{matriks Jacobian} ($J_F$) dengan metode \emph{Jacobi} pada SPL: keduanya berbeda.

\subsection{Penurunan Newton multivariabel}
Linearisasi Taylor orde pertama di titik $\mathbf{x}^{(k)}$ menghasilkan
\begin{equation}
\mathbf{F}(\mathbf{x}^{(k)}+\Delta\mathbf{x})\approx
\mathbf{F}(\mathbf{x}^{(k)})+J_F(\mathbf{x}^{(k)})\Delta\mathbf{x}.
\end{equation}
Kita meminta model linear ini sama dengan nol. Akibatnya, \textbf{setiap iterasi Newton menyelesaikan SPL}:
\begin{align}
J_F(\mathbf{x}^{(k)})\Delta\mathbf{x}^{(k)}&=-\mathbf{F}(\mathbf{x}^{(k)}),\\
\mathbf{x}^{(k+1)}&=\mathbf{x}^{(k)}+\Delta\mathbf{x}^{(k)}.
\end{align}
Di MATLAB gunakan \texttt{delta = -J(z)\textbackslash F(z)}, bukan menghitung \texttt{inv(J)}.

\subsection{Algoritma Newton multivariabel}
\begin{enumerate}
\item Tentukan $\mathbf{F}$, Jacobian $J_F$, tebakan awal $\mathbf{x}^{(0)}$, toleransi $\tau$, dan batas iterasi $N_{\max}$.
\item Pada iterasi $k$, hitung $\mathbf{F}_k=\mathbf{F}(\mathbf{x}^{(k)})$ dan $J_k=J_F(\mathbf{x}^{(k)})$.
\item Pastikan Jacobian tidak singular/nyaris singular; selesaikan SPL $J_k\Delta\mathbf{x}^{(k)}=-\mathbf{F}_k$.
\item Perbarui $\mathbf{x}^{(k+1)}=\mathbf{x}^{(k)}+\Delta\mathbf{x}^{(k)}$. Jika langkah penuh memperbesar residual, coba langkah teredam atau ubah tebakan awal.
\item Simpan galat perubahan relatif berskala dan residual:
\begin{equation}
 e_k=\max_i\frac{|x_i^{(k+1)}-x_i^{(k)}|}{\max(1,|x_i^{(k+1)}|)},
 \qquad r_k=\|\mathbf{F}(\mathbf{x}^{(k+1)})\|_\infty.
\end{equation}
\item Hentikan jika kriteria yang dipilih terpenuhi; laporkan \emph{kedua} ukuran dan status bila $N_{\max}$ tercapai.
\end{enumerate}
\Note{Tebakan awal adalah \emph{educated guess}, bukan harus angka tertentu. Patuhi domain model (misalnya konsentrasi positif), lihat grafik/arti variabel, dan hindari tebakan pada Jacobian singular. Newton lokal tidak menjamin konvergensi dari sebarang tebakan.}

\subsection{Contoh terhitung}
Tentukan penyelesaian
\begin{align}
f_1(x,y)&=x^2+y^2-5=0,\\
f_2(x,y)&=x-y-1=0.
\end{align}
Dengan $\mathbf{x}^{(0)}=(1.5,0.5)^T$, diperoleh
\begin{equation}
\mathbf{F}(\mathbf{x}^{(0)})=\begin{bmatrix}-2.5\\0\end{bmatrix},\qquad
J_F(\mathbf{x}^{(0)})=\begin{bmatrix}3&1\\1&-1\end{bmatrix}.
\end{equation}
Selesaikan $J_F\Delta\mathbf{x}=-\mathbf{F}$ untuk memperoleh $\Delta\mathbf{x}=(0.625,0.625)^T$; jadi $\mathbf{x}^{(1)}=(2.125,1.125)^T$. Lanjutkan sampai residual kecil. Periksa solusi analitik $(2,1)$ dan $(-1,-2)$: sistem ini mempunyai lebih dari satu akar.

\subsection{Pisahkan function dan driver}
\textbf{Function:} \texttt{newton\_spnl.m} memuat algoritma dan \texttt{history}. \textbf{Driver:} \texttt{driver\_spnl.m} menetapkan persamaan, Jacobian, tebakan awal, tabel, serta grafik.
\MatlabFile{newton\_spnl.m (function)}
\lstinputlisting[language=Matlab]{matlab/newton_spnl.m}
\MatlabFile{driver\_spnl.m (driver)}
\lstinputlisting[language=Matlab]{matlab/driver_spnl.m}

\subsection{Latihan SPNL}
\begin{enumerate}
\item Turunkan $J_F$ sendiri sebelum menjalankan kode; hitung satu iterasi manual.
\item Jalankan dari tiga tebakan awal dan identifikasi solusi yang diperoleh; perhatikan domain yang relevan.
\item Buat \texttt{semilogy} untuk \texttt{history.errs} dan \texttt{history.residual}; jelaskan jika suatu grafik stagnan.
\item Modifikasi driver untuk sistem lain, sedangkan function Newton tetap sama.
\end{enumerate}

\section{Bagian B --- Regresi linear dan nonlinear}
\subsection{Data dan definisi dua model}
Pada contoh kelas, gunakan \textbf{data yang sama} untuk kedua model:
\begin{center}
\begin{tabular}{c|ccccc}\hline
$x_i$&0.25&0.75&1.25&1.75&2.25\\
$y_i$&0.28&0.57&0.68&0.74&0.79\\\hline
\end{tabular}
\end{center}
Model linear: $\widehat y_i=a+b x_i$. Model nonlinear (saturasi sesuai contoh kelas):
\begin{equation}
\widehat y_i=f(x_i;\mathbf{a})=a_0\bigl(1-e^{-a_1x_i}\bigr),
\qquad \mathbf{a}=(a_0,a_1)^T,\quad \mathbf{a}^{(0)}=(1,1)^T.
\label{eq:model_saturasi}
\end{equation}
Model linear mempunyai koefisien linear dalam parameter. Model saturasi nonlinear dalam parameter $a_1$, sehingga membutuhkan iterasi untuk mendapatkan nilai least squares.

\subsection{Regresi linear: dari OLS sampai SPL}
Residual $r_i=y_i-(a+bx_i)$. \emph{Ordinary least squares} (OLS) meminimalkan
\begin{equation}
S(a,b)=\sum_{i=1}^n\bigl[y_i-a-bx_i\bigr]^2.
\end{equation}
Setel kedua turunan parsial objektif ke nol:
\begin{align}
\frac{\partial S}{\partial a}&=-2\sum_i(y_i-a-bx_i)=0,\\
\frac{\partial S}{\partial b}&=-2\sum_ix_i(y_i-a-bx_i)=0.
\end{align}
Diperoleh SPL \emph{normal equations}:
\begin{equation}
\underbrace{\begin{bmatrix}n&\sum x_i\\\sum x_i&\sum x_i^2\end{bmatrix}}_{X^TX}
\underbrace{\begin{bmatrix}a\\b\end{bmatrix}}_{\boldsymbol\beta}
=\underbrace{\begin{bmatrix}\sum y_i\\\sum x_i y_i\end{bmatrix}}_{X^T\mathbf{y}},\quad
X=\begin{bmatrix}1&x_1\\\vdots&\vdots\\1&x_n\end{bmatrix}.
\label{eq:linear_normal}
\end{equation}
Untuk data kelas,
\begin{equation}
\begin{bmatrix}5&6.25\\6.25&10.3125\end{bmatrix}
\begin{bmatrix}a\\b\end{bmatrix}
=\begin{bmatrix}3.06\\4.42\end{bmatrix},\qquad
\widehat{\boldsymbol\beta}=\begin{bmatrix}0.3145\\0.2380\end{bmatrix}.
\end{equation}
Dalam perangkat lunak, \texttt{X\textbackslash y} menyelesaikan OLS lebih stabil daripada \texttt{inv(X'*X)*X'*y}; persamaan normal dipakai untuk \textbf{memahami formulasi}, bukan perintah wajib menghitung invers.
\MatlabFile{regresi\_linear.m (function)}
\lstinputlisting[language=Matlab]{matlab/regresi_linear.m}

\subsection{Regresi nonlinear: fungsi, Jacobian, dan pembaruan parameter}
Untuk model \eqref{eq:model_saturasi}, residual dan objektifnya ialah
\begin{equation}
D_i(\mathbf{a})=y_i-f(x_i;\mathbf{a}),\qquad
S(\mathbf{a})=\tfrac12\sum_iD_i(\mathbf{a})^2.
\end{equation}
Turunan parsial \textbf{model terhadap parameter}:
\begin{align}
\frac{\partial f}{\partial a_0}&=1-e^{-a_1x},\\
\frac{\partial f}{\partial a_1}&=a_0x e^{-a_1x}.
\end{align}
Susun Jacobian model $Z\in\mathbb{R}^{n\times2}$ dengan baris
\begin{equation}
Z_{i,:}(\mathbf{a})=
\begin{bmatrix}1-e^{-a_1x_i}&a_0x_i e^{-a_1x_i}\end{bmatrix}.
\end{equation}
Di $\mathbf{a}^{(0)}=(1,1)^T$ berlaku (pembulatan empat desimal):
\begin{equation}
Z_0=\begin{bmatrix}
0.2212&0.1947\\0.5276&0.3543\\0.7135&0.3581\\0.8262&0.3041\\0.8946&0.2371
\end{bmatrix},\qquad
\mathbf{D}_0=\begin{bmatrix}0.0588\\0.0424\\-0.0335\\-0.0862\\-0.1046\end{bmatrix}.
\label{eq:class_Z_D}
\end{equation}
Ini adalah matriks dan residual pada \textbf{gambar contoh kelas}; nilai lengkap tanpa pembulatan digunakan dalam program.

\subsection{Menurunkan algoritma Gauss--Newton}
Linearisasi model $f(\mathbf{a}+\Delta\mathbf{a})\approx f(\mathbf{a})+Z\Delta\mathbf{a}$ menghasilkan
\begin{equation}
\mathbf{D}(\mathbf{a}+\Delta\mathbf{a})\approx\mathbf{D}-Z\Delta\mathbf{a}.
\end{equation}
Cari $\Delta\mathbf{a}$ yang meminimalkan $\|\mathbf{D}-Z\Delta\mathbf{a}\|_2^2$. Persamaan normal untuk langkah adalah
\begin{equation}
\underbrace{(Z^TZ)}_{2\times2}\Delta\mathbf{a}=Z^T\mathbf{D},
\label{eq:gn_normal}
\end{equation}
lalu perbarui parameter:
\begin{equation}
\mathbf{a}^{(k+1)}=\mathbf{a}^{(k)}+\alpha_k\Delta\mathbf{a}^{(k)},\qquad 0<\alpha_k\le1.
\label{eq:gn_update}
\end{equation}
Secara numerik, implementasi \texttt{delta = Z\textbackslash D} memecahkan masalah least squares tanpa membentuk $Z^TZ$ secara eksplisit. \texttt{alpha} dapat dibagi dua bila objektif justru membesar. Untuk data contoh, langkah penuh pertama sekitar $\Delta\mathbf{a}=(-0.27148,\,0.50193)^T$.

\subsection{Algoritma fungsi regresi nonlinear}
\begin{enumerate}
\item Input $x,y$, model $f(x;\mathbf{a})$, Jacobian model $Z(x,\mathbf{a})$, tebakan awal, toleransi dan batas iterasi.
\item Evaluasi vektor residual $\mathbf{D}=\mathbf{y}-\widehat{\mathbf{y}}$ dan Jacobian $Z$ berukuran $n\times p$.
\item Periksa dimensi dan rank, lalu selesaikan masalah least squares $Z\Delta\mathbf{a}\approx\mathbf{D}$ (ekuivalen dengan \eqref{eq:gn_normal} bila rank penuh).
\item Uji parameter $\mathbf{a}+\alpha\Delta\mathbf{a}$; kurangi $\alpha$ jika SSE membesar.
\item Catat parameter, \texttt{errs} dan SSE untuk grafik; ulangi sampai toleransi atau iterasi maksimum.
\item Periksa kewajaran parameter, pola residual dan kualitas hasil; perubahan parameter kecil saja tidak membuktikan kecocokan data baik.
\end{enumerate}
\MatlabFile{regresi\_nonlinear.m (function)}
\lstinputlisting[language=Matlab]{matlab/regresi_nonlinear.m}
\MatlabFile{driver\_regresi.m (driver)}
\lstinputlisting[language=Matlab]{matlab/driver_regresi.m}

\subsection{Evaluasi model dan grafik}
Dengan $r_i=y_i-\widehat y_i$, $\bar y=\frac1n\sum_i y_i$, dan $\mathrm{SSE}=\sum r_i^2$:
\begin{align}
\mathrm{MSE}&=\frac1n\sum_i r_i^2, &
\mathrm{RMSE}&=\sqrt{\mathrm{MSE}},\\
\mathrm{MAPE}&=\frac{100}{n}\sum_i\left|\frac{r_i}{y_i}\right|, &
R^2&=1-\frac{\mathrm{SSE}}{\sum_i(y_i-\bar y)^2}.
\end{align}
MAPE klasik tidak terdefinisi untuk $y_i=0$; $R^2$ tidak terdefinisi jika seluruh $y_i$ sama. Jangan menuliskan hasil metrik tanpa memeriksa kondisi tersebut. \texttt{driver\_regresi.m} menghasilkan grafik titik data dan kedua model, grafik residual, serta \texttt{semilogy} galat parameter per iterasi.

\section{Latihan terintegrasi dari Lembar Kerja}
\begin{enumerate}
\item[\textbf{B4.1 (L1)}] Susun $x^2+y^2-5=0$ dan $x-y-1=0$ dalam bentuk vektor $\mathbf{F}(x,y)=\mathbf{0}$. Jelaskan hubungan dengan pencarian akar skalar.
\item[\textbf{B4.2 (L2)}] Tentukan Jacobian $J_F(x,y)$ dan daerah di mana Jacobian tidak invertibel.
\item[\textbf{B4.3 (L3)}] Mulai dari $(1.5,0.5)$, kerjakan satu langkah Newton multivariabel secara manual, lalu cocokkan dengan \texttt{driver\_spnl.m}.
\item[\textbf{B4.4 (L4, di kelas)}] Sesuai Lembar Kerja, gunakan data $ (1,0.5),(2,2.5),(3,2.0),(4,4.0),(5,3.5)$. (a) Cari model regresi linear $\widehat y=a+bx$ dengan OLS, turunkan SPL normalnya. (b) \textbf{Tambahan nonlinear:} pada \emph{data yang sama} fit $\widehat y=a_0(1-e^{-a_1x})$, bentuk Jacobian parameter $Z$, lalu jalankan Gauss--Newton. (c) Hitung MSE, RMSE, MAPE, $R^2$ dan buat grafik data-model, residual serta \texttt{errs}. Bandingkan pula dengan contoh data kelas di bagian sebelumnya.
\item[\textbf{B4.5 (L5, tugas)}] Gunakan dua belas data saturasi pada tabel berikut. Model \textbf{default} adalah $\widehat y=a_0(1-e^{-a_1x})$. Hitung parameter dan keempat metrik serta grafik \texttt{errs}. Untuk tantangan, coba model saturasi rasional $\widehat y=c_0x/(c_1+x)$ pada \textbf{data yang sama}; turunkan Jacobian parameter dan bandingkan metriknya. Jelaskan perbedaan bentuk model dan batasan parameter. Untuk mempertahankan butir asli B4.5, tuliskan juga objektif $S(a,b)=\sum_i(y_i-ae^{bx_i})^2$ dan kondisi optimalitas $\partial_a S=\partial_b S=0$ untuk model eksponensial $ae^{bx}$.
\end{enumerate}
\begin{center}\small
\begin{tabular}{c|rrrrrr}\hline
$x$&0.25&0.75&1.25&1.75&2.25&2.75\\
$y$&0.165&0.440&0.612&0.771&0.859&0.945\\\hline
$x$&3.25&3.75&4.25&4.75&5.25&5.75\\
$y$&0.988&1.044&1.065&1.097&1.102&1.118\\\hline
\end{tabular}
\end{center}
\Note{Data tugas adalah data sintetis pembelajaran, bukan hasil pengukuran lapangan. Bila ingin mencoba model eksponensial tumbuh $\widehat y=ae^{bx}$, gunakan dataset tambahan dalam \texttt{matlab/data\_eksponensial.csv}; model itu \emph{tidak} mempunyai plateau dan karena itu tidak otomatis cocok untuk data saturasi. Untuk dua parameter $a,b$, turunan Jacobian-nya adalah $\partial f/\partial a=e^{bx}$ dan $\partial f/\partial b=ax e^{bx}$.}

\subsection*{Dataset opsional: pertumbuhan eksponensial (12 titik)}
Untuk mahasiswa yang ingin mencoba fungsi $\widehat y=ae^{bx}$ dari butir awal B4.5, gunakan data sintetis \emph{berbeda} berikut, bukan data saturasi di atas:
\begin{center}\small
\begin{tabular}{c|rrrrrr}\hline
$x$&0.5&1.0&1.5&2.0&2.5&3.0\\
$y$&0.806&0.920&1.159&1.410&1.641&2.057\\\hline
$x$&3.5&4.0&4.5&5.0&5.5&6.0\\
$y$&2.488&2.932&3.604&4.366&5.235&6.395\\\hline
\end{tabular}
\end{center}
Fit $f(x;a,b)=ae^{bx}$ dengan Jacobian model $Z_i=[e^{bx_i},\; ax_i e^{bx_i}]$. Sebagai tantangan lain, bandingkan dengan fungsi pangkat $f(x;c,d)=cx^d$ (untuk $x>0$), dengan turunan parameter $\partial_c f=x^d$ dan $\partial_d f=cx^d\ln x$. Hitung empat metrik pada \textbf{data yang sama} dan gambarkan \texttt{errs}. Jangan menggabungkan dataset pertumbuhan ini dengan dataset saturasi pada satu fitting.

\section{Tugas refleksi komputasi}
Jalankan seluruh driver dengan data sesuai instruksi. Jelaskan isi input/output tiap function, tunjukkan apa yang berubah ketika tebakan awal diganti, serta simpan grafik \texttt{errs} dan residual dengan label lengkap.
